\subsection*{Summary}

Existing literature illustrates that while the landscape for \ac{HPC} workloads has shifted significantly over time, their access patterns can still be accurately modeled by existing \ac{HPC} storage performance benchmarking tools: large sequential read behaviour of traditional \ac{HPC} applications can be emulated accurately by \ac{IOR}, while small-file-\acs{I/O} and other metadata-heavy operations in \ac{AI}/\ac{ML} can be modeled by mdtest. Ancillary workloads, such as common operations performed by users interactively on their home directories are similarly metadata-heavy, and are likewise modelable using mdtest. Benchmarking suites such as the IO500 leverage both tools in order to establish a bounding box for file system performance, capturing both best and worst-case performance metrics for both data throughput and metadata operations.

In addition, Ceph has evolved significantly from its early versions. In particular, the transition from FileStore to BlueStore for \ac{OSD} backing storage resolved many of Ceph's major performance bottlenecks, and has allowed Ceph to more effectively utilize modern technologies such as \ac{NVMe} flash storage. However, while reports have shown high aggregate throughput results, per-client performance may be constrained by communication overhead, possibly exacerbated by the \ac{CPU} overhead of the \ac{TCP/IP} networking stack on which Ceph relies. Furthermore, the compounding factor of underlying physical disk performance, particularly with \acs{HDD}-backed \acp{OSD}, obscure some of the intrinsic architectural limitations of Ceph's software stack. 

Given that Ceph is significantly more reliant on its software logic than incumbent solutions such as Lustre, it is necessary to decouple the limitations of the hardware from the storage system itself. Therefore, our methodology employs memory-backed storage to mitigate storage performance as a compounding factor, allowing us to directly stress and analyze the performance of Ceph's abstraction layers---in particular that of \ac{RADOS} and \ac{CephFS}.